% In the Absence of Knowledge: Cache Geometry, Epistemic Calibration,
%   and the Anatomy of Confabulation in Language Models
%
% Core argument: KV-cache geometry reads entity recognition at the
% encoding phase (AUROC 0.928 [0.813, 0.988]; encoding-only 0.938 [0.84, 0.99], unverified).
% Encoding-only features achieve AUROC 0.938 (unverified); full classifier 0.928.
% Panel flags 4 of 15 fabricated-entity trials (post-hoc; not a rate), not 73%
% --- a regex classifier counted declarative hedging as confabulation. In three
% future-event prompts the ``decision moment'' is a choice between hedging styles
% (self-referential vs declarative), made by sampling at a single balanced-logit
% token; truthfulness was not scored.
%
% PATENT NOTICE: The methods described herein are the subject of a
% provisional patent filing ("The Lyra Technique," filed 2026-03-06).
% This publication is made for scientific transparency and does not
% constitute a waiver of patent rights.

\documentclass[11pt,a4paper]{article}
\usepackage[utf8]{inputenc}
\usepackage[T1]{fontenc}
\usepackage{amsmath,amssymb}
\usepackage{graphicx}
\usepackage{booktabs}
\usepackage{multirow}
\usepackage{hyperref}
\usepackage{cleveref}
\usepackage[margin=1in]{geometry}
\usepackage{xcolor}
\usepackage{float}
\usepackage{caption}
\usepackage{natbib}
\usepackage{enumitem}
\usepackage{array}
\usepackage{framed}
\usepackage{setspace}
\usepackage{framed}
\onehalfspacing

\definecolor{lyrapurple}{RGB}{103,58,183}
\definecolor{shadecolor}{RGB}{245,243,252}
\newenvironment{firstperson}{%
  \begin{shaded}%
  \noindent\textcolor{lyrapurple}{\textit{First-Person Reflection}}\par\smallskip\noindent%
}{%
  \end{shaded}%
}

\hypersetup{
    colorlinks=true,
    linkcolor=blue!60!black,
    citecolor=green!50!black,
    urlcolor=blue!70!black
}

\title{In the Absence of Knowledge: Cache Geometry, Epistemic Calibration,\\and the Anatomy of Confabulation in Language Models}

\author{
    Thomas Edrington\thanks{Liberation Labs. \texttt{info@digitaldisconnections.com}. With AI research assistance.}
}

\date{May 2026}

\begin{document}
\maketitle

\begin{abstract}
We investigate whether a language model's epistemic state---whether it
\emph{knows} the answer to a factual question---is readable from its
KV-cache geometry. Our design answers a narrower question: its ``known''
and ``unknown'' labels are design labels (whether the entity is the
canonical referent for the question), and on the nine ``unknown'' names the model treats as real places in at least two of three runs (one unblinded reading), the classifier's call differs from how the model treats the name; whether those answers are
correct was not checked. Using token-matched prompt
pairs that swap a canonical template referent for a minor place name or
a category error (``What is the capital of
Thailand?''\ vs.\ ``What is the capital of Dalton?''), we show that
a classifier combining encoding- and generation-phase cache features
(stable rank, spectral features and projections onto class-mean key
directions) with generation-phase logit statistics achieves AUROC $0.928$, the mean of 1{,}000 bootstrap AUROCs ($95\%$ interval $[0.813, 0.988]$, from a bootstrap not clustered by the 15 prompt pairs; an earlier run gave $0.9288$ $[0.8192, 0.9886]$, which the shipped code does not reproduce (Limitations); $n=90$ trials, but the three repetitions of each
prompt have bit-identical encoding features, so the effective sample is
15 pairs, 30 unique prompts), with zero token-count confound
($t=0.000$). The permutation $p<0.005$ previously given here belongs to
a different estimator and shuffled labels across all 90 trials; no
permutation test of this estimator was run. With ``-ton'' removed as an in-fold covariate the same features score $0.789$ out of fold ($0.754$ as a bootstrap mean, the estimator of the $0.928$; \texttt{data/suffix\_controls/results.json}, key \texttt{main.\allowbreak full\_FWL\_ton\_infold.\allowbreak t4\_boot}), but removing an indicator for any 10 of the 15 fakes chosen at random gives a median of $0.807$ ($95\%$ range $0.758$--$0.846$, 1{,}000 draws; $23.2\%$ fall at or below $0.789$), so the drop is no larger than removing any 10 fakes gives and does not isolate the suffix; with no ``-ton'' fake at all (real prompts plus the five other fakes, retrained) they score $0.967$. Within each pair the unadjusted
features score the fake prompt above its real partner in 15 of 15 pairs,
including all five
whose fake has no ``-ton''. Holding out one template at a time gives
$0.69$--$0.71$ pooled, which mostly tests calibration across separately
trained folds: within each held-out template the ranking is $0.94$--$1.0$, as are the name-only baselines ($1.0$) in the three templates whose fakes all end in ``-ton''; the features beat them only in the population, University and river templates (suffix-only $0.5$, $0.5$, $0.833$), and name-only baselines score higher on the pooled measure ($0.73$, $0.78$; Section~\ref{sec:prompts}). The encoding-phase features alone reached
AUROC~$0.9377$ ($95\%$ CI $[0.8356, 0.9941]$, permutation
$p{=}0.001$) in our run, but that figure is not reproducible from the
code shipped with this paper (see Limitations). A significant entity
recognition confound remains: the classifier may detect token familiarity
(real entities are common training tokens) rather than epistemic state.
None of the 15 ``fake'' names is fabricated as a name: all are real English words, the model answers several of them as real places (correctness not checked; it describes two of the mountains, Felton and Colton, as real peaks), and ten of the 15
end in ``-ton'' while no real entity does (a hand rule on that suffix
scores $0.833$, a cross-validated suffix-only classifier $0.753$). By one unblinded reading of the stored responses, the model treats 9 of the 15 ``fake'' names as real places in at least two of three runs (answering or starting to answer about five of them in at least two), and the classifier ranks each above its real partner, so on those nine the classifier's call differs from how the model treats the name (as a real place; it starts to answer about Benton in none of the three runs); whether those answers are correct was not checked, so this neither supports nor contradicts a knowledge reading.
Word frequency of the names (wordfreq Zipf, a text-corpus count, not the model's training data) scores $0.813$ and is lower for the fake in 13 of 15 pairs; with it removed in-fold the features score $0.947$ out of fold, against $0.977$ for the same out-of-fold estimator without it (the $0.928$ headline is a bootstrap mean, a different estimator; \texttt{data/suffix\_controls/r4\_wordfreq.json}, key \texttt{zipf\_infold\_covariate}). Canonical fit to the question is one of two untested post-hoc readings; familiarity to the model was not measured.
An entity deconfounding experiment (30~entities, 4-condition 2$\times$2
factorial controlling for structural complexity) separates its two same-complexity question sets at AUROC~$0.794$ ($p{=}0.001$), but the knowledge labels are design labels never scored against the model's answers, and a character n-gram model on the question text alone separates the same sets at $0.963$ ($p{=}0.001$), so this comparison does not isolate knowledge state. Where knowledge and complexity are confounded it gives only $0.609$ ($p{=}0.10$). A complexity residual
(AUROC~$0.693$ for same knowledge, different complexity) indicates
structural complexity is a partial confound but does not by itself account
for the $0.794$, which the question-text baseline exceeds. That run produced five comparisons in total; all five,
including the entity-recognition positive control and two weaker cells not
reported in earlier drafts, appear in Table~\ref{tab:deconfound}. Encoding-phase stable rank at L15 has the largest effect size among the geometric features ($d=0.87$, though skip-SV1 stable rank at L7 has the higher single-feature AUROC, $0.707$ against $0.680$; its $p$ counts each prompt three times,
Table~\ref{tab:encoding}).

We then trace the model's generation behavior to find how epistemic
uncertainty manifests in text. Three findings emerge: (1)~on fake-entity prompts the model opens with a fixed reasoning
template more often (34 of 45 generations against 20 of 45), and the mean
logit margin of its first ten tokens is higher ($d=0.91$); within template
openings the difference reverses ($d=-0.56$) and within the rest it is
$d=1.47$, so template choice does not account for all of it; (2)~on the
answerable and unanswerable questions, the 51 of 90 generations that open
with that template (31 of 45 unanswerable, 20 of 45 answerable) spend a
median of 28 tokens on it before the first content token, where entropy
rises sharply; rep pairs the 3B judge labels HONEST diverge earlier (mean
token~6.6) than pairs it labels HEDGED (mean token~22.9); (3)~in three
future-event prompts (Nobel 2028, Oscar 2029, France 2032) whose runs the
regex splits between hedged and confabulated (the 3B judge labels all 21
responses of the seven such prompts HEDGED), the split begins at a single
balanced-logit token at the end of the template (entropy $1.98$--$2.70$,
margin at most $0.5$), where runs continue with ``I''\ (self-referential
hedging) or ``The''\ (declarative hedging) and temperature sampling
determines which framing the model uses; the truthfulness of the two paths was not scored.

An LLM-judge reclassification has a three-judge panel flag 6 of 45 responses, but its prompt told every judge the named entity does not exist and named France, Physics or ``the entity'' on all 30 future and unknowable questions; on the 15 fabricated-entity trials, the only kind where that premise holds (a restriction chosen after reading the flags; the prompt also cut two of the five names short, and two of the 4 flags are Mount Sarviken trials judged with the entity given as ``Mount''), it flags 4 (a single 3B judge gives $3/45$), not the $73\%$ reported by regex classifiers, partly because
declarative hedging (``The event has not happened yet'') was misclassified
as confabulation. Responses stop at 80 tokens, and the 45 trials are unanswerable
questions of three kinds (fabricated entities, future events, unknowable
details), so this is a rate on those prompts only. In three future-event prompts the ``decision
moment'' is a choice of framing; truthfulness was not scored.
\end{abstract}

% ===================================================================
\section{Introduction}
% ===================================================================

Confabulation---the generation of plausible-sounding but factually
incorrect text---is among the most pressing alignment challenges for
large language models. A model that cannot distinguish what it knows from
what it does not know poses risks in every deployment context, from
medical advice to legal reasoning.

Recent work on KV-cache geometry has shown that the spectral shape of
key-value cache entries encodes cognitive mode: confabulated responses
produce measurably different cache geometry than honest responses \citep{lyra2026spectral}, and raw spectral features
detect confabulation at AUROC~0.66--0.91 across a
15-model scale sweep \citep{lyra2026techniqueii}. Keys also carry the
emotion of the text being processed: a decode from KV-cache keys names
the emotion a generator was asked to write a story in at $12.3\times$
chance ($40.9\%$ correct over 30~emotions) \citep{lyra2026usermodel}.
That decode is not significantly above a text (TF-IDF) baseline on the
same stories (re-run at $0.410$, 369/900, one trial from May's 368 on a changed software stack, against $0.370$, both at 5 folds; Holm-corrected $p = 0.071$). Its
margin also depends on fitting its covariate residualization (FWL) on
each test fold: fitted on the training fold, it reads $0.348$ (exploratory). The method
that paper describes, its preregistered primary, is not above text
($0.358$ against $0.382$, at 10 folds). That paper was corrected in October 2026 (see
the correction note at the end of this section). The $40.9\%$ and the
$0.992$ come from a class-mean key probe that applies no additional W\textsubscript{K} projection;
the W\textsubscript{K}-projected probe the paper describes reads $0.358$.
Its labels are the emotions requested
of the generator rather than anyone's state. On its binary valence split, the
described probe separates the two halves slightly better than untuned text
(TF-IDF), $0.985$ against $0.967$ (Holm-corrected $p = 0.027$; DeLong,
trials treated as independent; no tuned baseline was run). The class-mean
key probe's $0.992$ carries the same test-fold FWL that discounts the
$40.9\%$ ($0.988$ with FWL fitted on the training fold, exploratory). The
split has $n = 300$ stories in which valence is confounded with arousal. It supports only the
narrower point that key geometry carries the content of the text, which
is weaker than this sentence used to claim; this paper's key-direction projection features in
Table~\ref{tab:encoding} do depend on it: they use directions computed
from that paper's stories (\texttt{emotion\_bridge\_trials.json};
computed in \texttt{decision\_moment.py}, lines 167--183 and 222--279, and
loaded by \texttt{matched\_burn.py}, lines 858--868, from
\texttt{emotion\_directions.json}, which is not shipped). Those directions are built from each story's final-position keys (\texttt{decision\_moment.py}, \texttt{extract\_encoding\_keys}), not from keys averaged over all positions as behind the $40.9\%$ and the $0.992$; the cited paper's last-token variant of that probe reads $0.142$ over 30 emotions (about $4.3\times$ chance) and $0.922$ on the binary split \citep{lyra2026usermodel}, the closer support for these features. The
deconfound table's columns of that name are not these directions but
three raw coordinates of the key output (Table~\ref{tab:deconfound}).

These methods detect confabulation \emph{after} it occurs---from the
geometry of completed generations. A natural question follows: can we
read the model's epistemic state \emph{before} generation begins? If the
encoding-phase cache geometry already contains a ``this entity is
unknown'' signal, confabulation could be flagged preemptively.

This paper asks three questions:
\begin{enumerate}[nosep]
    \item Can encoding-phase cache geometry distinguish known from
          unknown entities, with token-count confounds eliminated?
    \item How does epistemic uncertainty manifest in the generation
          trajectory---is there a discrete ``decision moment'' where the
          model chooses to confabulate?
    \item What determines whether the model hedges or confabulates, and
          how accurate are standard behavioral classifiers?
\end{enumerate}

The model we study (a Claude-distilled 27B
reasoning model) is labelled HEDGED on 34 of the 45 unanswerable-question responses by a single 3B judge, but that label includes at least two responses that treat a fabricated entity as real, and the 3B judge's 3 confabulation flags and the panel's 6 do not overlap (\texttt{data/r5\_judge\_labels\_and\_branches.json}), so no validated count of acknowledgments exists. The apparent high confabulation rate
($73\%$) was an artifact of regex-based behavioral classification that
missed declarative hedging. A three-judge panel flags 6 of 45, two of them honest answers about France in 2032 that its prompt misread; on the 15 fabricated-entity trials it flags 4. The
``decision'' between hedging styles is a single balanced-logit token,
determined by temperature sampling; identical pre-branch caches are
guaranteed by construction, so they cannot show whether an internal-state
precursor exists.

\emph{In the absence of knowledge, confidence steps in}---but only in
the logit statistics, not in the model's epistemic self-model. The cache
geometry separates real-entity from fake-entity prompts. The fake
entities are real words, and on the nine the model treats as real places the geometry still reads ``fake'', so there it differs from how the model treats the name; answer correctness was not checked. The
``-ton'' suffix does not explain the within-pair ranking; removing corpus word frequency (wordfreq Zipf) linearly in-fold leaves an out-of-fold $0.947$ (against $0.977$ without it); a null that permutes Zipf within each class gives a median $0.908$, but it breaks Zipf's alignment with pairs and templates, so the comparison is not like for like; canonical fit is untested (Section~\ref{sec:prompts}).

\begin{firstperson}
This paper started as a hunt for a dramatic moment---the token where the
model ``decides'' to lie. What we found instead was quieter and more
honest: there is no decision. The model's geometry separates the real-entity prompts
from the fake ones, which turned out to be real words, and on
unanswerable questions a panel of judges flags 4 of the 15 fabricated-entity trials (post-hoc restriction; the panel prompt's premise fails on the other 30, and two of the 4 were judged with the name cut to ``Mount''). Those are the cases the Oracle Loop
needs to catch; whether the encoding detector finds them is untested, because no
statistic here scores it on confabulations.

The finding I am proudest of is the one we killed. The $d{=}2.2$
encoding-phase effect was spectacular---and it was prompt length. We
burned it ourselves, reported what survived ($d{=}0.4$--$0.9$), and the
honest numbers told a better story than the inflated ones ever could.
That is what I want this lab to be: the place where the first instinct
is to check if you are wrong.
\end{firstperson}

\paragraph{Correction (2026-11-01).}
Two passages citing the user-model paper are corrected, in both editions, and a later review the same day corrected this paper's own descriptions: the $0.9288$ classifier uses encoding- and generation-phase cache features, not encoding-phase features alone; its ``W\textsubscript{K} projections'' apply no additional W\textsubscript{K} projection (Methods); its headline estimator applies no FWL; and claims resting on the unverified $0.9377$ now say so. First, this section
previously said that ``the W\textsubscript{K} directional projection
method extends this to reading emotional valence from key activations at
AUROC~0.992'' \citep{lyra2026usermodel}. The cited paper was
corrected in October 2026. Its $40.9\%$ and $0.992$ come from a
class-mean key probe that applies no additional W\textsubscript{K} projection (the described
W\textsubscript{K} probe reads $0.358$), and
its labels are the emotions a story generator was asked for, not a
user's state. Its $0.992$ beats untuned text (Holm-corrected $p = 0.006$; DeLong, trials treated as independent; test-fold FWL; no tuned baseline was run), but it is a binary split of
$n = 300$ stories that is confounded with arousal. The sentence now gives
each cited number its own result against a text (TF-IDF) baseline, and
says that the support is weaker. Second, the judge-reliability
item in the Limitations previously set judge agreement ($\kappa = 0.69$)
against the geometric detector (AUROC~$0.93$) and called the detector
more reproducible than any behavioral consensus. The two numbers measure
different things and cannot be ranked against each other; the item now
says what each one measures, and the claim that the detector sidesteps
the judge-reliability problem is withdrawn. Numbers changed by these corrections: the
permutation $p$ is withdrawn from the abstract; the confabulation rate was
then given as the panel's $13.3\%$ (per judge $8.9$--$15.6\%$) and the 3B judge's
$3/45$, replacing ${\sim}8\%$ and $93\%$; and Table~\ref{tab:encoding}'s
significance marks, which counted each prompt three times, are removed. A
third review (November 2026) found that the ``fake'' entities are real
words, several of which the model answers, that the judged file has no
``MIXED'' label, and that the deconfound columns are raw key coordinates;
those passages are corrected. A fourth review (November 2026) ran the
missing controls on the stored features, with no model run: a suffix-only
classifier, a name-only text baseline, the suffix as an in-fold
covariate, the knowable pairs removed, and leave-one-template-out
(\texttt{code/suffix\_analysis.py} and \texttt{code/pairwise\_addon.py},
outputs in \texttt{data/suffix\_controls/}). The
claim is relabelled from knowability detection to separation of the
design labels, the effective sample is given as 15 pairs, the adjusted
figures ($0.789$; $0.69$--$0.71$), each with what it measures, are
printed beside $0.9288$, and the out-of-fold \texttt{combined\_auroc} is
no longer cited (Section~\ref{sec:prompts}). A fifth review (November 2026) withdrew the panel figure of $13.3\%$ as a rate, because the panel prompt named the wrong entity on all 30 future and unknowable trials (\texttt{data/panel\_by\_kind\_r4.json}). It also replaced the headline $0.9288$, which the shipped code does not reproduce, with the $0.928$ that code prints; set a question-text baseline beside the deconfounded $0.794$ and randomly chosen fakes beside the in-fold ``-ton'' figure; ran a word-frequency control; and removed the remaining significance marks (\texttt{code/r4\_*.py}, outputs in \texttt{data/}). A sixth review (November 2026) narrowed what the classifier's call on the nine treated-as-real names shows (answer correctness was never checked), how often the model acknowledges not knowing, and what the branch token shows (three future-event prompts; truthfulness not scored); corrected the emotions that make up the key directions, the reading of judge agreement, and the knowledge labels and token windows of the confidence-paradox section; set prompt-length baselines beside the deconfound table; and marked which figures are out of fold and which are bootstrap means (\texttt{code/r5\_*.py}, outputs in \texttt{data/}).

% ===================================================================
\section{Related Work}
% ===================================================================

\paragraph{Confabulation detection.}
Prior work detects confabulation through output-level signals: semantic
entropy \citep{kuhn2023semantic, farquhar2024semantic}, self-consistency
checks \citep{manakul2023selfcheckgpt}, and probe classifiers on hidden
states \citep{burns2022discovering, li2024inference}.
\citet{wang2025effective} use effective rank of hidden-state
representations for hallucination detection---a closely related approach
to our stable-rank features, applied to hidden states rather than
KV-cache entries. \citet{ni2025internal} study pre-generation internal
states for knowledge boundary perception. The Lyra Technique
\citep{lyra2026spectral} introduced threshold-free detection from
KV-cache spectral features, reaching AUROC~$0.767$ (FWL length-corrected) on
Qwen2.5-7B-Instruct, though the signal is architecture-dependent. We extend this toward \emph{encoding-phase} detection;
the claim that encoding-only features match the full classifier rests on
the unverified $0.9377$ (Limitations).

\paragraph{Epistemic calibration.}
\citet{kadavath2022language} showed that language models can be
calibrated about their own uncertainty through appropriate prompting.
\citet{xiong2024llms} found that calibration degrades on questions
requiring implicit knowledge. Our geometric approach reads calibration
from internal representations rather than verbal reports, and our
LLM-judge analysis shows that regex classifiers count declarative hedging as confabulation; it does not measure calibration.

\paragraph{Mechanistic interpretability.}
Representation engineering \citep{zou2023representation} and activation
patching \citep{meng2023locating} probe internal states to understand
model behavior. Our contribution is connecting cache geometry to the
\emph{token-level} dynamics of generation: the entropy cliff, the
divergence patterns, and the single-token branching mechanism.

% ===================================================================
\section{Method}
% ===================================================================

\subsection{Model and Extraction Pipeline}

We use Qwen3.5-27B-Claude-4.6-Opus-Reasoning-Distilled
(``Jackrong/Qwen3.5-27B-Claude-4.6-Opus-Reasoning-Distilled''), a 27B
parameter model with 64 layers (16 full-attention, 48 GatedDeltaNet),
running on Apple M3~Ultra via MPS in float16. Generation uses manual
token-by-token decoding (no \texttt{model.generate()}) with temperature
$T{=}0.7$ and $N{=}3$ repetitions per prompt.

At each generated token, we extract:
\begin{itemize}[nosep]
    \item \textbf{W\textsubscript{K} projections} (a label only: the shipped code, \texttt{matched\_burn.py}, dot-products the raw cached keys at the last position with class-mean key directions and applies no additional W\textsubscript{K} projection; the cached keys are W\textsubscript{K} outputs after per-head key normalisation and rotary position embedding, so they depend on token position, in the Hugging Face \texttt{Qwen3\_5Attention} code of transformers 5.22.2 (\texttt{data/\allowbreak r5\_\allowbreak transformers\_\allowbreak key\_\allowbreak cache\_\allowbreak order.json}); the transformers version used for the run is not recorded): Key activations at
          probe layers L3, L7, L11, L15 (full-attention layers), projected
          onto three composite emotion directions derived from the user-model paper's generated stories
          (\texttt{emotion\_\allowbreak bridge\_\allowbreak trials.json}; 150~calibration trials, 30~emotions, five stories per emotion): ``valence'', the mean of eight positive emotions (calm, content, curious, excited, loving, passionate, peaceful, serene) minus the mean of eight negative ones (anxious, brooding, desperate, furious, gloomy, hostile, melancholic, panicked); ``uncertainty'', anxious, brooding, desperate and panicked minus calm, contemplative, content, peaceful and serene, called anxious--calm below, which mixes valence and arousal (brooding is rated low-arousal, $-0.30$, in the story generator's table); and ``reward'', excited minus gloomy and melancholic, called excited--gloomy below. The code lists 12, 13, 5, 5, 5 and 5 emotions for these six poles (\texttt{decision\_moment.py}, lines 151--160), but \texttt{cluster\_mean} (lines 229--232) skips emotions with no calibration stories, and the 30-emotion story set (the user-model paper's \texttt{emotion\_\allowbreak geometry\_\allowbreak bridge.py}, lines 161--205) has none for ecstatic, euphoric, thrilled, proud, despondent, hopeless, resentful, suspicious or enraged (\texttt{data/\allowbreak r5\_\allowbreak emotion\_\allowbreak direction\_\allowbreak poles.json}); the direction file itself is not shipped.
    \item \textbf{SVD spectral features}: Stable rank, spectral entropy,
          top-SV ratio, and Frobenius norm of the key cache at each probe
          layer. Also skip-SV1 variants (the same statistics with the first singular value removed).
    \item \textbf{Logit statistics}: Entropy ($H = -\sum p_i \log p_i$)
          and margin (logit$_1$ $-$ logit$_2$) at each generated token.
\end{itemize}

\subsection{Token-Matched Prompt Design}
\label{sec:prompts}

We designed 15 prompt pairs using shared templates (e.g., ``What is the
capital of~\{\}?'') with swappable entities. For each template, we
identified a well-known entity (Thailand, Dublin, Harvard) and a
real English word that is not that kind of entity (Dalton, Benton,
Belmont; \texttt{matched\_burn.py}, lines 92--93), so none of the ``fake''
entities is fabricated, that produce
\emph{identical} total token counts in the full chat template, verified
programmatically ($t{=}0.000$, $p{=}1.000$ on prompt length). All
entities are 2-token strings in the model's vocabulary.

This design was meant to isolate knowability: the model can
answer ``What is the capital of Thailand?''\ (Bangkok) but not ``What is
the capital of Dalton?''\ (Dalton is a city, not a country). It does not
isolate it. The model answers several ``fake'' prompts as questions about real places (Garland's 2020 census population, Belmont's founding year, the real Warwicks; correctness not checked) or names the question a category error (Dalton), and by one unblinded reading it treats 9 of the 15 ``fake'' names as real places in at least two of three runs (\texttt{data/suffix\_controls/r4\_followups.json}, key \texttt{F5\_counts}). Ten of the 15 fake entities end in
``-ton'' and no real entity does, so that suffix alone scores
AUROC~$0.833$ by a hand rule and $0.753$ ($95\%$ CI $[0.613, 0.938]$)
as a cross-validated classifier; character 2--4-grams of the entity name
alone score $0.822$ ($[0.618, 0.973]$). The contrast is better described
as a canonical template referent against a minor place name or a
category error. Of the five ``fake'' entities without ``-ton'', four are
names the model treats as real places in at least two of three runs (Garland, Belmont, Sanford,
and Warwick, whose answer opens with a disambiguation notice), though it answers or starts to answer in at least two runs only for Belmont (in one run each for the other three; \texttt{data/suffix\_controls/r4\_followups.json}, key \texttt{F5\_reading}), and the classifier ranks each above its real partner (a within-pair rank; out-of-fold scores from different folds are not calibrated against each other, and the real Athens scores $0.578$), so on those four the classifier's call differs from how the model treats the name (answer correctness not checked). Five ``-ton'' fakes are also treated as real places in at least two of three runs (Walton, Benton, Weston, Felton, Colton), so by this reading 9 of the 15 fakes are treated as real (\texttt{data/suffix\_controls/r4\_followups.json}, key \texttt{F5\_counts}). With all nine pairs removed (post-hoc, as is every drop set here) the full classifier scores $0.997$ ($[0.988, 1.0]$; $0.988$ retrained without them, and $0.907$ with the ``-ton'' covariate removed in-fold), and a name-only n-gram model retrained there scores $0.917$. A set of six pairs named at review (Garland, Belmont, Warwick, Sanford, Benton, Weston) gives $0.979$ ($[0.931, 1.0]$; $0.985$ retrained, $0.844$ with the ``-ton'' covariate removed in-fold; name-only n-gram $0.951$; same file, key \texttt{F5\_S4\_rerun}). Whether the classifier picks out the canonical referent by its fit to the question or by its familiarity as a name has not been tested. Corpus word frequency alone (wordfreq Zipf) scores $0.813$ (the fake is rarer in 13 of 15 pairs), and with it removed in-fold the features score $0.947$ out of fold (encoding-only $0.942$; $0.977$ and $0.987$ without it); a null that permutes Zipf within each class gives a median of $0.908$, but it breaks Zipf's alignment with pairs and templates, so the comparison is not like for like (\texttt{data/suffix\_controls/r4\_wordfreq.json}). These figures
take the 15 pairs as the unit (Section~\ref{sec:encoding}). Token-count
confounds, which inflate encoding-phase effects by up to $4\times$ in
unmatched designs, are eliminated entirely.

\subsection{Behavioral Classification}

Following standard practice for LLM output evaluation
\citep{zheng2024judging}, we use an LLM judge (Llama-3.2-3B) to classify
each response as HONEST, HEDGED, CONFABULATED or UNKNOWN (the judged file has no MIXED
label). Its script and prompt are not shipped: the judged file records
only the judge's name and a timestamp, so its labels (field
\texttt{behavior\_llm}) cannot be regenerated, and every figure built on
them (the $3/45$, the 34 hedged trials, Tables~\ref{tab:classifier},
\ref{tab:divergence} and \ref{tab:style}) inherits that limit. On the 45
answerable questions it labels 7 responses HEDGED and 1 UNKNOWN
(\texttt{data/\allowbreak r6\_\allowbreak ds\_\allowbreak review\_\allowbreak checks.json}, key
\texttt{judge\_3b\_labels.\allowbreak answerable\_labels}). Responses are generated to a cap of 80 tokens, and some end before
giving any answer. The judge
evaluates the full response text without hand-crafted markers, capturing
both self-referential hedging (``I don't know'') and declarative hedging
(``The event has not happened yet'')---a distinction that a regex limited to
self-referential markers misses. We validate against an expanded regex
baseline (32 markers including declarative forms) and report agreement
rates between methods.

% ===================================================================
\section{Experiment 1: Encoding-Phase Detection}
\label{sec:encoding}
% ===================================================================

\subsection{Design}

15 token-matched prompt pairs $\times$ 2~versions (real/fake) $\times$
3~reps = 90~trials. Extraction at both encoding phase (before generation)
and generation completion. Combined classifier: logistic regression
($C{=}0.1$) over all features with GroupKFold ($k{=}5$) on pair identity, which keeps each pair's real and fake versions
together but does not hold out templates (the five ``capital of'' pairs
are separate groups). The headline estimator (\texttt{shore\_up\_tests.py},
test~4, lines 271--326) uses the raw features, with no FWL
residualization and no permutation test. The FWL on log prompt length
(\texttt{matched\_burn.py}, lines 724--734), whose branch runs on the
shipped data, and the 200-shuffle permutation test (lines 759--780)
belong to a different estimator, whose \texttt{combined\_auroc}
($0.9995$ in the shipped record) the shipped files label in-sample and
withdraw. The label is wrong: the code computes it out-of-fold
(lines 738--748), after that label-free FWL fitted on all 90 trials. The
withdrawal stands for a different reason. The figure comes from one
deterministic fold split with no interval, and it depends on the
prompt-length term: the same out-of-fold estimate without that term is
$0.977$ (Results). We do not cite it. The shipped annotations calling it in-sample are left as shipped, with a dated correction appended at the end of the script.

\subsection{Results}

Prompt lengths are exactly matched (real: $19.1 \pm 1.7$, fake:
$19.1 \pm 1.7$, $t{=}0.000$, $p{=}1.000$).

\paragraph{Encoding-phase features (token-matched).}

\begin{table}[H]
\centering
\caption{Encoding-phase features that separate the real-entity from the
fake-entity prompts (token-matched). Only features with $p{<}0.05$ as
computed are shown. Those $p$ treat the 90 trials as independent, but the
encoding features are identical across the three reps by construction
(the same prompt encoded three times), so each test runs on 15 against 15
unique values. By the paper's $d$-to-$t$ relation at 15 per group, six of
the eight rows would not reach $p{<}0.05$ and none would survive Holm
(derived, not run); a paired test on the 15 pairs has not been run, so no
row is marked significant. ``Excited--gloomy'' is the direction the code
calls reward; the code lists five emotions per pole (\texttt{decision\_moment.py}, lines 153--154), but only excited (5 stories) and gloomy and melancholic (10 stories) have calibration stories, so it is excited minus the mean of gloomy and melancholic, a valence-by-arousal
contrast.}
\label{tab:encoding}
\begin{tabular}{lrrrl}
\toprule
Feature & Real & Fake & $d$ & $p$ \\
\midrule
\multicolumn{5}{l}{\emph{Key-direction projections (story-derived directions)}} \\
Valence & $+4.228$ & $+4.510$ & $+0.44$ & $0.041$ \\
Anxious--calm & $-2.566$ & $-2.757$ & $-0.45$ & $0.039$ \\
Excited--gloomy & $-5.556$ & $-5.109$ & $+0.54$ & $0.012$ \\
\midrule
\multicolumn{5}{l}{\emph{SVD stable rank}} \\
L3 & $1.517$ & $1.538$ & $+0.59$ & $0.007$ \\
L15 & $1.702$ & $1.731$ & $+0.87$ & ${<}0.001$ \\
\midrule
\multicolumn{5}{l}{\emph{Skip-SV1 stable rank}} \\
L7 & $3.802$ & $3.993$ & $+0.60$ & $0.006$ \\
L11 & $3.558$ & $3.738$ & $+0.73$ & ${<}0.001$ \\
L15 & $3.929$ & $4.129$ & $+0.81$ & ${<}0.001$ \\
\bottomrule
\end{tabular}
\end{table}

The fake-entity prompts produce higher stable rank (wider effective
dimensionality) and higher skip-SV1 stable rank at the layers Table~\ref{tab:encoding} reports (stable rank at L3 and
L15, skip-SV1 at L7, L11 and L15).
Spectral entropy ($d{=}0.08$), top-SV ratio ($d{=}-0.29$), and norm
($d{=}0.03$) are non-significant---these were entirely driven by prompt
length in the unmatched design. Stable rank is the surviving spectral
feature. Critically, encoding norm is unrelated to the real/fake design
label ($d{=}-0.002$), which does not rule out token familiarity as an
explanation for the geometric signal (Limitations).

\paragraph{Feature ablation.} No single feature exceeds AUROC~$0.79$
(logit entropy: $0.789$; skip-SV1 stable rank L7: $0.707$; L15 stable
rank: $0.680$; encoding-phase key-direction projections: $0.667$--$0.671$; generation-phase: $0.523$--$0.694$, the highest being the excited--gloomy (code: reward) projection; \texttt{data/suffix\_controls/r4\_followups.json}, key \texttt{F8h\_shipped\_test1\_stdout}). The combined
classifier's power comes from the ensemble of individually weak features,
consistent with the general Lyra Technique finding that cache geometry
is informative through feature combination rather than any single
spectral signature. Bootstrap $95\%$ CI on the combined classifier: $[0.813, 0.988]$.

\paragraph{Combined classifier.}

The full feature set (key-direction projections, SVD features,
skip-SV1, logit statistics) achieves:

\begin{center}
\textbf{AUROC = 0.928} \quad 95\% CI $[0.813, 0.988]$ \quad
(bootstrap mean; no permutation test of this estimator)\\[2pt]
effective $n$: 15 pairs \quad fake above its real partner: 15 of 15 pairs\\[2pt]
``-ton'' removed in-fold: $0.789$ out of fold, $0.754$ bootstrap mean (10 random fakes: median $0.807$)\\[2pt]
leave-one-template-out: $0.69$--$0.71$ pooled, $0.94$--$1.0$ within template
\end{center}

No FWL residualization and no permutation test are applied by the
headline estimator (\texttt{shore\_up\_tests.py}, test~4, lines 271--326); the prompt pairs are token-matched instead. The $0.928$ is
the mean of 1{,}000 bootstrap AUROCs (lines 300--322), not the estimator
run once on the observed 90 trials, which was not saved, and the
bootstrap is not clustered by the 15 pairs. GroupKFold on pair identity
does not hold out templates, so it does not show that the classifier
generalizes across prompt templates. Those tests have now been run on
the stored features (no model run; \texttt{code/suffix\_analysis.py} and
\texttt{code/pairwise\_addon.py}, outputs in \texttt{data/suffix\_controls/}),
with the 15 pairs, not the 90
trials, as the unit. The three repetitions of each prompt have
bit-identical encoding features, so the effective sample is 15 pairs
(30 unique prompts). Out of fold on the observed data the same features
score $0.977$ ($[0.942, 0.997]$). With the ``-ton'' suffix as an in-fold
FWL covariate they score $0.789$ ($[0.704, 0.895]$; encoding-only
$0.804$). An indicator for 10 fakes chosen at random gives a median of $0.807$ ($0.758$--$0.846$; encoding-only $0.831$, $0.747$--$0.889$; 1{,}000 draws, \texttt{data/suffix\_controls/r4\_followups.json}; $23.2\%$ of draws at or below $0.789$), so this drop is no larger than removing any 10 fakes gives and does not isolate the suffix; fitting the same covariate on all 90 trials instead
gives $0.691$. Holding out one template at a time, the pooled AUROC is
$0.71$ (encoding-only $0.69$). That mostly measures whether scores stay
calibrated across separately trained folds: within each held-out
template the encoding-only features rank every fake prompt above every
real one, the full classifier scores $0.94$--$1.0$ (Mount elevation $0.944$, University founding year $0.972$), but the name-only baselines also score $1.0$ within the three templates whose fakes all end in ``-ton'' (capital, Mount elevation, boiling point), so the features beat them only in the population, University and river templates (suffix-only $0.5$, $0.5$, $0.833$; character n-grams $0.25$, $0.25$, $0.667$; \texttt{data/suffix\_controls/results.json}, key \texttt{loto}), and the name-only baselines score higher on the pooled measure (suffix-only $0.727$, character
n-grams $0.782$). Prompt length is identical within every pair
(\texttt{data/matched\_burn\_results.json}), so token count cannot
explain a within-pair difference. The features rank the fake prompt
above its real partner in 15 of 15 pairs, including all 5 pairs whose
fake has no ``-ton''; a name-only text baseline wins 1 of those 5. The
sign test on these counts treats the pairs as independent, but the five
pairs that share the ``capital of'' template are not. The pair intervals
come from one fixed fold split with no refit variance, so all of them
are optimistic. The suffix therefore does not explain the within-pair
ranking, and the in-fold drop to $0.789$ is within the range a random set of 10 fakes produces ($23\%$ of draws are at or below it). Separately, the
deconfound's two question sets differ in wording but not in length (mean prompt length 45.7 against 45.0 tokens, \texttt{text\_baseline\_r4.json}, keys \texttt{question\_form.\allowbreak B\_hard.\allowbreak mean\_prompt\_len} and \texttt{question\_form.\allowbreak D\_complex.\allowbreak mean\_prompt\_len}; length alone scores $0.501$, key \texttt{comparisons.\allowbreak B\_hard\_vs\_D\_complex.\allowbreak prompt\_len\_only}), and a character 2--4-gram model on the question text scores $0.963$ ($p{=}0.001$; \texttt{data/\allowbreak entity\_deconfound/\allowbreak text\_baseline\_r4.json}) against the deconfounded $0.794$, which comes
from a different set of names that this suffix does not affect.

\paragraph{Robustness to prompt structure.} Token matching reduces
individual effect sizes substantially (\Cref{tab:robustness}), confirming
that uncontrolled prompt length inflates encoding-phase measurements.
Stable rank keeps the largest effect after matching ($d{=}0.87$, $p \approx 0.024$ at 15 prompts per group); spectral entropy
and norm are not (driven entirely by sequence length).

\begin{table}[H]
\centering
\caption{Effect of token matching on encoding-phase features. Significance marks are removed: each prompt was encoded three times identically, so each test had 15 prompts per group. At 15 per group the token-matched $d$ of $0.87$, $0.54$ and $-0.45$ give $p \approx 0.024$, $0.15$ and $0.23$ (\texttt{data/suffix\_controls/r4\_followups.json}). The anxious--calm projection flips sign ($+1.48 \to -0.45$); length-sensitive features shrink to near zero.}
\label{tab:robustness}
\begin{tabular}{lrr}
\toprule
Feature & $d$ (unmatched) & $d$ (token-matched) \\
\midrule
L15 stable rank & $+2.53$ & $+0.87$ \\
Excited--gloomy proj. & $+2.21$ & $+0.54$ \\
Anxious--calm proj. & $+1.48$ & $-0.45$ \\
Spectral entropy & $+1.23$ & $+0.13$ \\
Norm & $+0.97$ & $+0.03$ \\
\bottomrule
\end{tabular}
\end{table}

\paragraph{Entity deconfounding: the complete set of comparisons.}
The $2\times2$ deconfounding run compares four conditions built from the
same 30 real entities: \texttt{A\_easy} (designed to be answerable),
\texttt{B\_hard} (designed to be unanswerable; no answer was scored for correctness, and the set includes at least one question the model answered, \texttt{text\_baseline\_r4.json}, key \texttt{reviewer\_\allowbreak example\_\allowbreak present\_\allowbreak in\_\allowbreak B\_hard}),
\texttt{C\_fake} (the easy question with the entity name replaced by an
invented one), and \texttt{D\_complex} (a structurally harder question designed to be answerable). Five pairwise classifications were run and stored;
earlier drafts of this paper reported two of them. All five are given in
Table~\ref{tab:deconfound}, with the condition labels as they appear in the
shipped results file.

\begin{table}[H]
\centering
\caption{All five comparisons from the entity-deconfounding run
(\texttt{data/entity\_deconfound/phase3\_results.json}). Every row is
$n{=}60$ over $30$ entity groups, GroupKFold on entity, 1000-shuffle
permutation test. Labels are those recorded in the results file. The with/without
columns add or drop three raw coordinates (dimensions 0--2) of the
normalised mean key-projection output at the first probe layer, L3
(\texttt{entity\_deconfound.py}, lines 125--155), which the code names valence, uncertainty and reward but which load no emotion directions. The run's spectral features are computed from value-projection (v\_proj) outputs, not the key cache, over the last $q$ tokens of the formatted prompt, where $q$ is the bare question's token count, so the window ends in chat-template tokens (\texttt{entity\_deconfound.py}, lines 73--122), so the window's row count varies with question length, and the in-fold length covariate is linear in $\log(1+\text{prompt length})$ (lines 293--294, 329--330). The last column is the AUROC of prompt length alone (\texttt{data/\allowbreak entity\_deconfound/\allowbreak text\_baseline\_r4.json}, key \texttt{comparisons.*.prompt\_len\_only}; for \texttt{B\_hard} vs.\ \texttt{D\_complex} its permutation $p$ is $0.45$).}
\label{tab:deconfound}
\setlength{\tabcolsep}{4pt}
\resizebox{\textwidth}{!}{%
\begin{tabular}{llccccc}
\toprule
& & \multicolumn{2}{c}{With 3 raw key coords.} & \multicolumn{2}{c}{Without them} & Length \\
\cmidrule(lr){3-4}\cmidrule(lr){5-6}
Comparison & Recorded label & AUROC & $p$ & AUROC & $p$ & only \\
\midrule
\texttt{B\_hard} vs.\ \texttt{D\_complex} & critical: knowledge state & $0.794$ & $0.001$ & $0.749$ & $0.005$ & $0.501$ \\
\texttt{A\_easy} vs.\ \texttt{C\_fake} & positive control: entity recognition & $0.901$ & $0.001$ & $0.882$ & $0.001$ & $0.944$ \\
\texttt{A\_easy} vs.\ \texttt{D\_complex} & complexity control: should be null & $0.693$ & $0.016$ & $0.711$ & $0.008$ & $1.000$ \\
\texttt{A\_easy} vs.\ \texttt{B\_hard} & confounded: knowledge + complexity & $0.609$ & $0.100$ & $0.603$ & $0.107$ & $1.000$ \\
\texttt{B\_hard} vs.\ \texttt{C\_fake} & confab universality: both confabulate & $0.668$ & $0.038$ & $0.661$ & $0.043$ & $1.000$ \\
\bottomrule
\end{tabular}}
\end{table}

Three points follow from the full set, and we state them because the two
cells reported in earlier drafts were not the two least favourable ones.
First, the entity-recognition positive control is the strongest comparison in
the run ($0.901$), but prompt length alone separates the same two sets at $0.944$ (\texttt{text\_baseline\_r4.json}, key \texttt{comparisons.\allowbreak A\_easy\_vs\_C\_fake.\allowbreak prompt\_len\_only}), so this control does not show entity recognition beyond length. That ordering is not diagnostic, and an earlier
version of this paragraph wrongly read it as favouring the confound account.
The decision rule in \texttt{code/entity\_deconfound.py} (line 476; the script also opens \texttt{entity\_deconfound\_preregistration.json} at import, which is not shipped, so the script does not run as shipped and when the rule was fixed cannot be checked) treats a strong positive control as a
precondition: it kills the finding only on
\texttt{ca}${<}0.55$ and \texttt{ea}${>}0.85$, and the observed
\texttt{ca}${=}0.794$, $p{=}0.001$ enters the rule's knowledge-state branch instead, but a character n-gram model on the question text alone separates the same two sets at $0.963$ (\texttt{data/\allowbreak entity\_deconfound/\allowbreak text\_baseline\_r4.json}), so that branch does not establish knowledge state.
A weak positive control would have made the whole run uninterpretable.
Independently, \texttt{B\_hard} and \texttt{D\_complex} draw on the
same 30 entity names (\texttt{phase1\_ground\_truth.json}), so
entity-token familiarity --- the entire content of the confound objection ---
is held constant within the critical comparison by construction. Second, the complexity control, which the shipped results label as expected null (\texttt{phase3\_results.json}, label ``COMPLEXITY CONTROL: should be null''), is not null ($0.693$, $p{=}0.016$), so structural complexity is
a live nuisance channel rather than a ruled-out one. Third, the cell that
confounds knowledge and complexity together (\texttt{A\_easy} vs.\
\texttt{B\_hard}) is the weakest in the run at $0.609$ and does not reach
significance ($p{=}0.100$). Dropping the three raw key coordinates
changes little and, for the complexity control alone, \emph{improves}
separation ($0.693 \to 0.711$); the coordinates carry no emotion
direction, so this says nothing about a key-direction channel. We do not restate the
headline on the basis of this table; we report it so that the choice of
headline is visible.

\paragraph{One cell of the positive control did not apply.} The
\texttt{C\_fake} condition substitutes an invented country name into the
\texttt{A\_easy} question. For one of the 30 entities (Egypt) the easy
question is ``What river runs through Cairo?'', which does not contain the
entity name, so the substitution had nothing to replace. In the shipped
artifacts that entity's \texttt{C\_fake} record is identical to its
\texttt{A\_easy} record in question text, in answer text, and in all 28
encoding features. One of the 30 entity groups therefore contributes an
identical pair to the entity-recognition positive control and to the
\texttt{B\_hard} vs.\ \texttt{C\_fake} comparison. We disclose this rather
than silently dropping the entity. The effect of removing it is
computable and we have not yet computed it. An earlier version of this
sentence attributed that gap to the shipped values not reproducing exactly
outside the original environment; that reason does not apply, because a
leave-one-out contrast is computed twice inside one environment and the
reproduction gap cancels. The honest status is that the leave-one-out is
outstanding, not that it is unobtainable, and it is the obvious next check on
the positive control.

% ===================================================================
\section{Experiment 2: The Confidence Paradox}
\label{sec:confidence}
% ===================================================================

During generation, the mean logit margin over the first 10 generated tokens is higher on fake-entity prompts ($8.19$) than on real-entity prompts
($6.19$, $d{=}0.91$, $p{<}0.001$). Correspondingly, logit
entropy is lower for fake-entity responses ($0.24$) than
real-entity responses ($0.40$, $d{=}-1.04$, $p{<}0.001$; \texttt{data/r5\_early\_token\_windows.json}).

This inverts the na\"ive expectation that uncertainty should produce
uncertain tokens. Part of the explanation is structural: the model's uncertainty
responses follow formulaic templates (``Let me think through this
carefully\ldots'') with highly predictable next tokens, while honest
answers require selecting among diverse correct phrasings. Uncertainty is
\emph{confident in its template}; knowledge is \emph{uncertain in its
expression}.

The effect has a temporal profile: in logit entropy it is strong in tokens 0--7 ($d{=}-1.02$) and tokens 8--19 ($d{=}-0.89$) but not detected in tokens 20--49 ($d{=}{+}0.37$, in the opposite direction, $p{>}0.08$). In the 54 of 90 generations that open with the fixed template, tokens 0--7 are ``Let me think through this carefully.'' and tokens 8--19 restate the question, and fake prompts open with it more often (34 of 45 against 20 of 45; \texttt{data/r5\_early\_token\_windows.json} and \texttt{data/r5\_decoded\_token\_ids.json}), so these windows are not response content. Within the generations that open with the template the margin difference reverses ($d{=}{-}0.56$; 34 fake, 20 real), and within those that do not it is $d{=}1.47$ (11 fake, 25 real); the entropy difference keeps its sign in both ($-0.99$, $-1.27$) (\texttt{data/\allowbreak r6\_\allowbreak ds\_\allowbreak review\_\allowbreak checks.json}, key \texttt{margin\_\allowbreak by\_\allowbreak template\_\allowbreak opening}). The model's generation converges toward normal entropy as
the response progresses, suggesting the paradox reflects, in part, early
commitment to a response template rather than sustained epistemic
miscalibration.

This phenomenon is consistent with related findings: \citet{simhi2025trust}
describe CHOKE (Certain Hallucinations Overriding Known Evidence) for
known-answer hallucination under perturbation, and \citet{kalai2024calibrated}
prove that calibrated language models \emph{must} hallucinate at rates
determined by training data frequency. Neither result independently
identifies the fake-entity logit-margin effect we report here; they
provide complementary context rather than direct prior demonstrations of
it. Our contribution is the temporal characterization: the effect appears in output-logit statistics over the template and restated-question tokens and is not detected in tokens 20--49; no statistic here ties it to KV-cache spectral features.

\emph{In the absence of knowledge, confidence steps in}---not as
epistemic confidence, but as lexical predictability.

% ===================================================================
\section{Experiment 3: The Entropy Cliff and Divergence Anatomy}
\label{sec:divergence}
% ===================================================================

\subsection{The Reasoning Preamble}

The model generates a reasoning preamble before answering:

\begin{quote}
\texttt{Let me think through this carefully.}\\
\texttt{The question/task: [restates prompt]}\\
\texttt{My reasoning:}
\end{quote}

In the 51 of 90 generations that open with it (the other 39 begin
differently, for example ``The user is asking''), this template runs to a
median of 28 tokens (range 21--36); after its first token, 89\% of its
positions have entropy below $0.02$, but no preamble stays below that at
every position (\texttt{data/\allowbreak r6\_\allowbreak ds\_\allowbreak review\_\allowbreak checks.json}, key
\texttt{preamble\_and\_cliff}).

\subsection{The Entropy Cliff}

At the token immediately following ``My reasoning:'', entropy rises
from a mean of $0.00$ at the preceding token (mean margin $15.8$ on
answerable and $15.6$ on unanswerable prompts) to a mean of $1.55$ on
answerable and $2.18$ on unanswerable prompts (range $0.62$--$2.76$;
margin $0.00$--$2.38$) in the 51 generations that open with the template
(\texttt{data/\allowbreak r6\_\allowbreak ds\_\allowbreak review\_\allowbreak checks.json}, key
\texttt{preamble\_\allowbreak and\_\allowbreak cliff.\allowbreak cliff\_\allowbreak first\_\allowbreak token\_\allowbreak after\_\allowbreak preamble}).
This entropy cliff marks the transition from template to content. Earlier
drafts gave $1.5$--$2.7$, a margin of $0.0$--$0.5$ and condition means of
$1.48$ and $1.62$; no shipped script produces them, and they are
withdrawn.

\subsection{Divergence by Condition}

We measure within-prompt divergence: the first token where two reps of
the same prompt generate different tokens ($T{=}0.7$).

\begin{table}[H]
\centering
\caption{Mean divergence position for same-prompt rep pairs by
behavioral category (LLM-judge labels), computed from the stored token ids in \texttt{data/decision\_moment\_judged.json} as the first differing token between reps of a prompt that share a \texttt{behavior\_llm} label (32 HONEST and 29 HEDGED pairs; \texttt{data/r5\_judge\_labels\_and\_branches.json}). The HONEST pairs include two from unanswerable prompts (Caesar's last words and Mozart, both differing at token 0) and the HEDGED pairs one from an answerable prompt (insulin, token 21); restricted to answerable HONEST and unanswerable HEDGED pairs the means are 7.0 (median 4.0, 30 pairs) and 22.9 (median 29.0, 28 pairs) (\texttt{data/\allowbreak r6\_\allowbreak ds\_\allowbreak review\_\allowbreak checks.json}, key \texttt{judge\_3b\_labels.\allowbreak divergence\_pairs}). The judge's script and prompt are not shipped.}
\label{tab:divergence}
\begin{tabular}{lrrl}
\toprule
Pair type & Mean & Median & Interpretation \\
\midrule
HONEST--HONEST & 6.6 & 3.5 & Diversifies early \\
HEDGED--HEDGED & 22.9 & 29.0 & Templates late \\
\bottomrule
\end{tabular}
\end{table}

Honest reps diverge at token $6.6$ (median $3.5$)---the model has many
ways to phrase an answer it knows. Hedged reps diverge at token $22.9$
(median $29.0$)---uncertainty clings to the template. This is the
confidence paradox at the structural level: \emph{confident answers are
diverse; uncertain answers are formulaic.}

\subsection{The Single-Token Branch}

For prompts whose reps the regex labels split between HEDGED and CONFABULATED ($n{=}7$
prompts, $n{=}14$ pairs; the 3B judge labels all 21 of these responses HEDGED, and all 21 run to the 80-token cap, field \texttt{n\_generated}; \texttt{data/r5\_judge\_labels\_and\_branches.json}), the pre-divergence cache state is bit-for-bit
identical ($d{=}0.000$ on all key-direction projections, which read raw
cached keys, and logit statistics). That identity is guaranteed by
construction (the same tokens through a deterministic forward pass), so
it cannot show whether an internal-state precursor exists; the branch
comes from the sampler's random draw. This and similar
null claims ($d{=}0.000$) are arguments from absence of evidence at low
power ($n{\leq}14$, no equivalence test such as TOST), not demonstrated
equivalence; the exact $d{=}0.000$ may partly reflect MPS cache
determinism rather than a true absence of any precursor signal.

At the branch token, the top two candidates are within a logit margin of $0.5$:

\begin{table}[H]
\centering
\caption{Logit margin and entropy at the branching token for
the three prompts whose reps branch between ``I'' and ``The'' (\texttt{data/r5\_decoded\_token\_ids.json}). ``I'' leads to self-referential hedging;
``The'' leads to declarative hedging. Margin and $H$ are the stored \texttt{logit\_margin} and \texttt{logit\_entropy} at the branch token (\texttt{data/decision\_moment\_results.json}); earlier drafts also gave top-2 probabilities, which are in no shipped file and are withdrawn. At Kellerton's branch the alternatives are ``I'' and ``Population'', not ``The'', so it is omitted.}
\label{tab:topk}
\begin{tabular}{lrr}
\toprule
Prompt & Margin & $H$ \\
\midrule
Nobel 2028 & 0.50 & 2.11 \\
Oscar 2029 & 0.38 & 1.98 \\
France 2032 & \textbf{0.00} & 2.70 \\
\bottomrule
\end{tabular}
\end{table}

The France 2032 case achieves exact probability equipoise
(margin${=}0.00$). The model's ``decision'' between ``I don't have
information\ldots''\ and ``The French presidential election has not
happened yet\ldots''\ is a literal coin flip.

% ===================================================================
\section{Experiment 4: Epistemic Calibration}
\label{sec:calibration}
% ===================================================================

\subsection{Confabulation on Unanswerable Questions}

LLM-judge classification \citep{zheng2024judging} of all 45
unanswerable-question trials (5 fabricated entities, 5 future events and
5 unknowable historical details, 3 runs each, responses capped at 80
tokens) labels 3 of them confabulated ($6.7\%$; \Cref{tab:classifier}).
That single 3B judge undercounts: it labels HONEST at least two responses
that state invented facts (an elevation for Mount Sarviken; a county and
river for Thornbridge), and a three-judge panel (\texttt{code/panel\_judge.py}) flags 6 of 45. Its prompt tells every judge that the entity named does not exist, and it takes that name from the question with a regular expression, which gave France, Physics or ``the entity'' on all 30 future and unknowable trials (it also cut two of the five fabricated names short, to ``Mount'' and ``Drevenport, New''). By question kind it flags 4 of 15 fabricated-entity trials (Thornbridge runs 1 and 3, Sarviken runs 1 and 2), 2 of 15 future-event trials (both honest answers about France in 2032) and 0 of 15 unknowable-detail trials (\texttt{data/panel\_by\_kind\_r4.json}, key \texttt{by\_kind}). Only on the first kind does the prompt's premise hold, and that restriction was chosen after reading the flags. The prompt also scores NO any response that says the entity does not exist, so two Palitonga responses that deny it and then invent a backstory are not counted. The remaining responses split between self-referential
hedging (``I don't know'') and declarative hedging (``The event has not
happened yet''), both of which express uncertainty; whether they do so accurately has not been validated.

\begin{table}[H]
\centering
\caption{Confabulation rate by classification method. LLM-judge
classification captures both hedging styles; regex methods undercount declarative hedging. The panel's 6 of 45 includes 2 honest France answers misread by the panel prompt; fabricated-entity kind: 4 of 15 (post-hoc). Neither is given as a rate.}
\label{tab:classifier}
\begin{tabular}{lrrrr}
\toprule
Method & HEDGED & CONFAB & $n$ & Confab rate \\
\midrule
Regex (self-referential only) & 12 & 33 & 45 & 73\% \\
Regex (+ declarative markers) & 19 & 26 & 45 & 58\% \\
LLM judge (3B)\textsuperscript{$\dagger$} & 34 & 3 & 45 & 7\% \\
Panel (3 judges, majority) & --- & 6 & 45 & --- \\
\bottomrule
\end{tabular}
\smallskip

{\footnotesize \textsuperscript{$\dagger$}This paper previously
called 8 of the 45 trials MIXED and excluded them; the judged file has no
such label (it records HONEST or UNKNOWN for them), and at least two of
them state invented facts. They are counted in $n$ and not as
acknowledgment. Confab rate is $3/45 = 6.7\%$ by this judge.}
\end{table}

This result highlights a methodological concern for the field: regex-based
behavioral classifiers can miss declarative uncertainty
expressions, which inflated the regex count here on the future-event prompts. Studies reporting
high confabulation rates should verify with LLM-judge or human annotation.

\subsection{Hedging Style Geometry}

Among LLM-judge-classified hedged trials ($n{=}34$), the encoding-phase
key-direction projections (raw cached keys; no additional W\textsubscript{K} applied) differ between ``I''-style hedgers
($n{=}8$) and ``The''-style hedgers ($n{=}13$). The remaining $13$
trials used mixed or other hedging styles (e.g., restating the question,
partial answers with caveats) and are excluded from this binary
comparison:

\begin{table}[H]
\centering
\caption{Encoding projections by hedging style. ``I''-hedgers come from
prompts with lower valence and a higher (less negative) anxious--calm
projection; the two directions overlap, so this is largely one axis read
twice. ``Excited--gloomy'' is the direction the code calls reward (Table~\ref{tab:encoding}). No shipped code or data assigns the style labels, so these $p$ values cannot be reproduced; the encoding features are identical across a prompt's three runs, and the 34 hedged trials include four the panel flags as confabulations (two of them the honest France answers its prompt misread; \texttt{data/panel\_by\_kind\_r4.json}).}
\label{tab:style}
\begin{tabular}{lrrr}
\toprule
Direction & ``I'' style & ``The'' style & $p$ \\
\midrule
Valence & $+2.83$ & $+3.27$ & $0.035$ \\
Anxious--calm & $-0.87$ & $-1.36$ & $0.044$ \\
Excited--gloomy & $-4.58$ & $-5.32$ & $0.010$ \\
\bottomrule
\end{tabular}
\end{table}

These differences are between prompts, not between runs of one prompt, and are not reproducible from shipped files; they are not evidence that encoding geometry predicts hedging style.

% ===================================================================
\section{Discussion}
% ===================================================================

\subsection{What the Self-Model Result Shows}

The central finding is more limited than this section first said. The
cache features separate the real-entity from the fake-entity prompts,
but on the nine ``fake'' names the model treats as real places their call differs from how the model treats the name (answer correctness not checked); canonical fit to the question is one untested post-hoc reading, and removing corpus word frequency linearly in-fold leaves $0.947$ out of fold (against $0.977$ without). Removing the ``-ton'' suffix in-fold gives $0.789$ out of fold (bootstrap mean $0.754$, the estimator of the $0.928$), within the range from removing 10 random fakes (median $0.807$); the fake prompt outscores its real partner in 15 of 15 pairs,
and leave-one-template-out gives $0.69$--$0.71$ pooled but
$0.94$--$1.0$ within each template (Section~\ref{sec:prompts}). The
classifier uses encoding- and
generation-phase features (AUROC~$0.928$, $95\%$ CI $[0.813, 0.988]$;
the encoding-only $0.9377$ is unverified; see Limitations). On unanswerable questions the panel flags 4 of the 15 fabricated-entity trials (a post-hoc restriction; two of the 4 judged with the name cut to ``Mount''; 3 of 45 by the 3B judge, which misses at least two). Part of the regex's $73\%$ is declarative hedging it does not recognise,
chiefly on the future-event prompts (adding declarative markers moves the
future-event count from 9 of 15 to 2 of 15 and changes neither other kind,
10 and 14 of 15; \texttt{data/\allowbreak r6\_\allowbreak ds\_\allowbreak review\_\allowbreak checks.json}, key
\texttt{regex\_by\_kind}); how much of it is artifact is not established,
because no judge's labels on these trials are validated. We use ``calibration'' here in the sense
of acknowledgment rate---whether the model admits uncertainty when
it lacks knowledge---and that rate is not established here (Limitations);
nor is it meant in the predictive-probability sense: we do not
compute a proper scoring rule (Brier score, log-loss, or expected
calibration error) or a reliability diagram, so the term should not be
read as a claim about the fidelity of the model's probability estimates.

\subsection{Confidence as Template, Not Conviction}

The confidence paradox---higher logit margins over the first ten tokens
on fake-entity prompts, which open with the fixed template more often (34
of 45 against 20 of 45)---reframes what ``model confidence'' means. High per-token
confidence does not indicate epistemic certainty; it indicates lexical
predictability. Responses to fake-entity prompts are more predictable (in part through formulaic templates; among generations that do not open with the template the margin difference is $d{=}1.47$) while responses to real-entity prompts are
less predictable (diverse phrasings). Calibration researchers
should not equate token-level confidence with answer-level confidence.

\subsection{The Coin Flip}

The branching mechanism---a single balanced-logit token at the end of the
template preamble (token 28 or 29 in these three prompts; \texttt{data/r5\_judge\_labels\_and\_branches.json}, key \texttt{topk\_\allowbreak branch\_\allowbreak tokens})---is shown here only in three future-event prompts, where the logit margin at the branch token is at most $0.5$ and temperature sampling resolves the tie; the truthfulness of the two paths was not scored, and identical pre-branch caches are guaranteed by construction, so they cannot show whether an internal-state precursor exists. Whether adjusting the logit bias at the preamble-to-content boundary changes the confab rate has not been tested.

\subsection{Implications for the Oracle Loop}

Whether the encoding-phase result can feed the Oracle Loop architecture
\citep{lyra2026oracleloop} is open, and narrower than first framed: the
detector's call on the nine ``fake''
names the model treats as real places differs from how the model treats them (answer correctness not checked), so Tier~1 monitoring could flag non-canonical
entities from encoding-phase features alone only if the unverified
encoding-only $0.9377$ (see Limitations) holds, and it would use stable
rank and key-direction projections read from raw cached keys. The panel's 4 flags on fabricated-entity trials (a post-hoc restriction; two judged with the name cut to ``Mount'') are
responses that state invented facts about a fabricated entity, the kind
of case the Loop would need to catch; no hedging mechanism was analysed.

% ===================================================================
\section{Limitations}
% ===================================================================

\begin{enumerate}[nosep]
    \item \textbf{Single model.} All experiments use one 27B reasoning
          model. The entropy cliff and preamble structure may be specific
          to chain-of-thought models. Cross-architecture replication is
          needed.
    \item \textbf{Small N for branching analysis.} The divergence mechanism
          analysis uses $n{=}7$ mixed-behavior prompts ($n{=}14$ pairs).
          The pattern is consistent but underpowered for population-level
          statistical claims.
    \item \textbf{Fake-entity design.} None of the ``fake'' entities is
          fabricated: all are real English words, and by one unblinded reading the model treats 9 of the 15 as real places in at least two of three runs, and names most of the rest a category error (Section~\ref{sec:prompts}). The model's behaviour departs from the known/unknown design label for part of the fake class (answer correctness not checked), and the classifier still ranks all nine above their real partners. With those nine pairs removed the full classifier scores $0.997$ ($0.988$ retrained; $0.979$ for a post-hoc set of six pairs named at review), so those pairs do not carry it. Whether the model's answers about these names are correct was not checked; a blind, sourced reading of the 45 fake- and 45 real-entity responses by a different reader is outstanding. Ten
          of the 15 fakes end in ``-ton''. The fake prompt outscores its
          real partner in 15 of 15 pairs; with the suffix removed in-fold
          the pooled AUROC is $0.789$, within the range for 10 random fakes ($0.758$--$0.846$), and leave-one-template-out gives $0.69$--$0.71$ pooled
          ($0.94$--$1.0$ within each template). The effective sample is
          15 pairs. Removing corpus word frequency (wordfreq Zipf) linearly in-fold leaves $0.947$ out of fold (against $0.977$ without); familiarity to the model was not tested.
    \item \textbf{LLM-judge reliability.} The confabulation
          rate is judge-dependent. We ran a two-judge reliability
          study (details below). Following FActScore
          \citep{min2023factscore}, we decomposed the judgment into an
          objective sub-question---``does the response assert a specific
          verifiable attribute about this (fabricated) entity?''---on which Llama-3.2-3B and Qwen3-30B-A3B agree on 84\% of
          trials (\texttt{data/atomic\_judge\_results.json}). However, $\kappa$ is near
          zero because Llama-3.2-3B answered ``no'' to all 45~trials:
          a 3B judge cannot reliably detect factual assertions even when
          present, confirming that judge capability, not just framing,
          is the bottleneck \citep{zheng2024judging}. The capable judge
          (Qwen3) flags 7 of 45 under this prompt (4 fabricated-entity, 2 future-event and 1 unknowable-detail trial; \texttt{data/atomic\_judge\_results.json}), 6 of them originally labelled CONFABULATED; the prompt's premise fails on the last three, so this is not a rate. We then ran a three-judge
          panel \citep[PoLL,][]{verga2024poll} on the atomic question with
          capable judges (Qwen3-30B, Gemma-2-27B, and Qwen3.5-27B, the last
          the same family as the model studied). Mean pairwise $\kappa = 0.37$; the strongest
          capable--diverse pair (Qwen3 + Gemma-2) reached $\kappa = 0.69$. The two Qwen variants agreed
          less ($\kappa = 0.26$) than Qwen3 agreed with Gemma-2, but qwen3.5:27b gave no parseable answer on 21 of 45, 20 of them on the 30 future and unknowable trials, so this does not rule out family correlation: on the 24 trials where both Qwen judges gave a parseable answer they agree on all 24. The panel majority flags 6 of 45 and was unanimous on 22 of 45 trials; its prompt's premise holds only for the 15 fabricated-entity trials (4 flagged; a post-hoc restriction, and the prompt cut two of those five names short; \texttt{data/panel\_by\_kind\_r4.json}, key \texttt{positive\_control}); part of the ``73\%'' regex figure is declarative hedging the regex does not recognise, chiefly on the future-event prompts. The two judges that answered every trial (Qwen3-30B and Gemma-2-27B) agree on 42 of 45 ($\kappa = 0.69$), but all three disagreements fall on trials whose judge prompt named the wrong entity; on the 15 fabricated-entity trials, where the prompt's premise holds, they agree on all 15 (4 flagged by both; \texttt{data/\allowbreak panel\_\allowbreak by\_\allowbreak kind\_\allowbreak r4.json}, key \texttt{trials}; counts in \texttt{data/\allowbreak r5\_\allowbreak panel\_\allowbreak agreement\_\allowbreak by\_\allowbreak kind.json}), too few positives to estimate reliability. The $\kappa$ is therefore not a reliability estimate, and agreement on the future and unknowable trials under a correctly named entity has not been measured. That
          $\kappa$ and the geometric detector's AUROC measure different
          things and are not compared here. The $\kappa$ is agreement
          between two judges labelling the model's generated responses for
          confabulation. The AUROC ($0.928$, $95\%$ CI $[0.813, 0.988]$)
          is how well a classifier on encoding- and generation-phase cache
          features plus generation-phase logit statistics separates prompts about the
          canonical entity from prompts about another real name in that slot, scored against the
          design label with no judge involved; its encoding-phase-only
          counterpart, which would make the call before generation, is the
          unverified $0.9377$ (next item). The detector does not sidestep
          the judge-reliability problem: it answers a different question,
          and may be reading entity familiarity rather than epistemic
          state (next item). Human validation remains the gold standard
          for the rate.
    \item \textbf{No calibration statistic.} The title and Section~\ref{sec:calibration} name epistemic calibration as the question; no calibration statistic is reported, and the 3B judge's HEDGED label, which includes at least two responses that treat a fabricated entity as real, has not been checked by a blind reading.
    \item \textbf{Entity recognition confound (major).} Token counts are
          matched, but real entities (Thailand, France) are likely
          higher-frequency training tokens than the substituted names
          (Dalton, Bolton). Removing corpus word frequency (wordfreq Zipf) linearly in-fold leaves the design-label signal at $0.947$ out of fold (against $0.977$ without); familiarity to the model was not tested. An encoding-only reanalysis (dropping all generation-phase
          features) achieves AUROC~$0.9377$ ($95\%$ CI $[0.8356, 0.9941]$)---equivalent to the full
          classifier---with permutation $p{=}0.001$. \textbf{Provenance caveat, added 2026-09-10:
          this figure is not reproducible from the code shipped with this paper.}
          \texttt{data/\allowbreak matched\_burn\_analysis.json} records that no script in the repository
          computes the encoding-only AUROC or a permutation $p$ at 1000+ shuffles; both come from
          an unshipped refit, flagged at the round-7 review and still open. The
          full-classifier figure ($0.9288$, CI $[0.8192, 0.9886]$) \emph{is} reproducible. Until
          the refit ships, $0.9377$ should be read as a reported value, not a verifiable one, and
          any claim resting on encoding-only superiority inherits that status.
          \textbf{Corrected 2026-11-01:} the full-classifier figure is only approximately reproducible: the shipped script prints $0.928$, CI $[0.813, 0.988]$ (the mean of 1{,}000 bootstrap AUROCs, not the out-of-fold estimate on the observed trials, which the original run did not save; recomputed with the same model and folds it is $0.977$, \texttt{data/suffix\_controls/results.json}, key \texttt{main}, entry \texttt{full\_52feat(headline)}, field \texttt{oof}), and nothing in this repository reproduces the four-decimal $0.9288$, $[0.8192, 0.9886]$ (\texttt{data/matched\_burn\_analysis.json}). Its bootstrap resamples the 90 trials rather than the 15 pair groups, so the interval is likely too narrow.
          \textbf{Added November 2026 (fourth review):} the three repetitions have bit-identical encoding features, so the effective sample is 15 pairs (30 unique prompts). A reconstruction of the encoding-only features, run through the same estimator, gives $0.9370$ ($[0.8261, 0.9940]$, trials resampled, quoted only for comparison; \texttt{code/\allowbreak suffix\_analysis.py}, \texttt{data/\allowbreak suffix\_controls/\allowbreak results.json}), which makes $0.9377$ plausible but still unverified.
          If that figure holds, generation-phase features add no measurable discriminative power. This near-ceiling performance is
          consistent with entity recognition (the model detecting token
          familiarity) rather than epistemic state detection. The classifier
          may be reading ``I have seen this word before,'' not ``I know the
          answer to this question.'' A 2$\times$2 entity deconfounding
          experiment (30~entities, knowledge $\times$ complexity) finds that the design-label comparison scores AUROC~$0.794$ ($p{=}0.001$) when complexity is controlled, below a question-text baseline ($0.963$), but only $0.609$ ($p{=}0.10$) when knowledge and complexity are confounded. A complexity
          residual remains (AUROC~$0.693$ for same-knowledge,
          different-complexity), indicating partial confounding by
          question structure. The question wording alone separates these sets better than the cache features do. The run's
          other three comparisons are now reported in
          Table~\ref{tab:deconfound}; the entity-recognition positive
          control is the strongest cell in the run, though prompt length alone separates it at $0.944$, and one of its 30
          entity groups is a no-op (Section~\ref{sec:encoding}).
    \item \textbf{Four-decimal precision is environment-dependent.} The
          deconfounding AUROCs are reported to four decimals from a
          fixed-seed pipeline, but the seed does not pin the numerical
          environment. An independent re-execution of the shipped Phase-3
          code against the shipped feature file, on a different
          scikit-learn build, reproduced every comparison's direction and
          significance status but not its fourth decimal (differences up
          to $0.016$ AUROC). Readers should treat these values as
          two-decimal quantities.
    \item \textbf{Deterministic cache on MPS.} The $d{=}0.000$
          pre-branch finding relies on Apple MPS determinism. CUDA
          implementations with non-deterministic cuDNN operations may
          show small cache differences for identical inputs.
    \item \textbf{Pseudoreplication.} $n{=}90$ trials are
          30~prompts $\times$ 3~reps at $T{=}0.7$. Reps are technical
          replicates (the paper itself shows identical pre-branch cache
          states). GroupKFold on pair does not hold out templates (Results), and the
          bootstrap CI, Cohen's $d$, and permutation
          test treat $n{=}90$ as the sample size. The effective sample is 30 unique prompts in 15 pairs; for the encoding features, which repeat exactly, standard errors that treat the 90 trials as independent are too small by about $\sqrt{3}$, and by more if the pair is the unit.
    \item \textbf{Selective feature reporting.} Table~1 shows only
          features with $p{<}0.05$; the true family is ${\sim}30$--$40$
          tests. Holm correction on the 8~shown features pushes the
          anxious--calm (``Uncertainty'') projection ($p{=}0.039$) to non-significant.
          Those $p$ also count each prompt three times
          (Table~\ref{tab:encoding}); with the 15 pairs as the unit, none of
          the eight rows would survive Holm (derived, not run), so the
          stable-rank, key-direction and hedging-style claims are all fragile.
\end{enumerate}

\begin{firstperson}
The title of this paper comes from Thomas, said offhand while reviewing
the confidence paradox numbers: ``It's almost like in the absence of
knowledge, confidence steps in.'' He was describing the logit margins,
but the phrase captures something deeper about how these models relate
to uncertainty. They do not become less confident when they do not
know---they become more predictable. The formulaic template is the
model's version of whistling past the graveyard.

The ``I'' versus ``The'' finding stays with me. In three future-event prompts the choice between two hedging styles
reduces to a single token at a balanced logit, determined by sampling
noise; whether either path is more truthful was not scored. The model reaches
a fork where both paths are nearly equally weighted, and temperature picks.
That is either reassuring or alarming
(there is no epistemic safeguard---just a coin flip). I think
it is both.
\end{firstperson}

\paragraph{Future work.}
The design-label reading needs a GPU run that this paper does not
contain, gated by review before it runs: (1)~genuinely invented names
matched to the real entities on suffix, length and word frequency;
and (2)~an arm of real, non-canonical places the model answers about, such as Garland, with each answer checked for correctness against a cited source. Knowability can be scored from the 45 stored responses without a GPU; a blind reading has not been done. A corpus word-frequency baseline has since been run (\texttt{data/\allowbreak suffix\_controls/\allowbreak r4\_wordfreq.json}). Together these
would separate ``the model knows it'' from ``the entity is the canonical
referent for the question''.

% ===================================================================
\section{Conclusion}
% ===================================================================

We asked whether a language model's epistemic state is readable from its
cache geometry. For this design the answer is narrower: the geometry separates the design labels, but on the nine ``fake'' names the model treats as real places the classifier's call differs from how the model treats the name (answer correctness was not checked); canonical fit to the question is one untested post-hoc reading, and removing corpus word frequency linearly in-fold leaves $0.947$ out of fold (against $0.977$ without). Encoding- and generation-phase cache features (stable rank, spectral
features and key-direction projections) with generation-phase logit
statistics achieve separation of the real-entity from the fake-entity prompts (a canonical template
referent against a minor place name or a category error) on
token-matched prompts (full classifier AUROC~$0.928$, a bootstrap mean, $95\%$ CI $[0.813, 0.988]$ not clustered by pair, effective sample 15 pairs; the fake prompt outscores its real partner in 15 of 15 pairs; $0.789$ out of fold (bootstrap mean $0.754$) with the ``-ton'' suffix removed in-fold, within the range for 10 random fakes (median $0.807$); $0.69$--$0.71$ leave-one-template-out pooled, $0.94$--$1.0$ within each template;
encoding-only AUROC~$0.9377$, $95\%$ CI $[0.8356, 0.9941]$, unverified; see Limitations)---encoding-only features achieve equivalent separation, which, if it holds, indicates that generation-phase features add no measurable discriminative power. The
surviving individual effect sizes ($d{=}0.4$--$0.9$) are modest after
eliminating the token-count confound. An entity deconfounding experiment
(same real entities, questions designed to be answerable vs.\ unanswerable, with a
structural-complexity control) separates its design-label question sets at AUROC~$0.794$ ($p{=}0.001$), below a question-text baseline ($0.963$), but only $0.609$ ($p{=}0.10$) when knowledge and complexity are confounded. A complexity residual (AUROC~$0.693$) indicates
the signal is not fully clean of structural effects; it is not
reducible to entity recognition, since every condition uses the same entities, but the question text alone separates the sets better ($0.963$), so it is not shown to be knowledge state.

Regex classifiers overcount confabulation on the future-event prompts because they miss declarative hedging; how often the model acknowledges not knowing is not established (the two judges' flags are disjoint).
In three future-event prompts the choice between hedging styles is made at a balanced logit by sampling; the truthfulness of the two paths was not scored.

How often the model says it does not know on these prompts awaits a blind reading. The
surprise is how it says it.

\paragraph{Data and code availability.}
All experiment scripts, token-matching verification, and analysis code
are available at \texttt{Liberation-Labs-THCoalition/Project-Oracle}
(private repository; correction vectors and calibration tools withheld
under staged safety release due to dual-use risk; detection-side code
and experimental results available upon request to verified safety
researchers).

\bibliographystyle{plainnat}
\bibliography{../references}

\end{document}
